% !TeX root = ../navier-stokes-manuscript.tex
% ============================================================
% 3.1 The Terminal Quantifier
% ============================================================
\subsection{The Terminal Quantifier}\label{sub:TheTerminalQuantifier}

Fix again the rapid-decay datum \(u_0\), the viscosity \(\nu>0\), and their
single maximal normalized history \((u,p)\).  Suppose its maximal time is
finite and write \(I_*=(0,T_*)\).  Every coordinate below is a derivative of
this history.  The datum belongs to the smooth solenoidal class
\begin{equation}\label{eq:SectionThreeSmoothSolenoidalClass}
  \Hinfty_\sigma(\R^3)
  :=
  \left\{
    v\in\bigcap_{m\in\Nzero}H^m(\R^3;\R^3)
    :\nabla\cdot v=0
  \right\}.
\end{equation}

Spatial decay fixes pressure from the velocity without an additive constant.
With \(R_j=\partial_j(-\Delta)^{-1/2}\),
\begin{equation}\label{eq:SectionThreePressureNormalization}
  -\Delta p=\partial_i\partial_j(u_i u_j),
  \qquad
  p=R_iR_j(u_i u_j).
\end{equation}
Direct time differentiation, with \(V_n=\partial_t^nu\) and
\(Q_n=\partial_t^np\), gives
\begin{align}
  Q_n
  &=R_iR_j\sum_{a=0}^{n}\binom{n}{a}V_{a,i}V_{n-a,j},
  \label{eq:SectionThreePressureRecurrence}\\
  V_{n+1}
  &=\nu\Delta V_n
  -\sum_{a=0}^{n}\binom{n}{a}(V_a\cdot\nabla)V_{n-a}
  -\nabla Q_n.
  \label{eq:SectionThreeVelocityRecurrence}
\end{align}
Every division of the time derivatives between the two transport factors remains
in the binomial sum.  Thus each temporal row still contains the full transport,
pressure, viscosity, and divergence-free evolution.

Choose finite temporal depth \(N\), spatial height \(M\), and a strict cutoff
\(T<T_*\).  The adjacent-row norm is
\begin{equation}\label{eq:SectionThreeRectangleNorm}
\begin{aligned}
  \mathfrak R_{N,M}(u;T)
  :={}&
  \sum_{n=0}^{N}
  \norm{\partial_t^n u}_{L^2(0,T;H^{M+1})}^{2}
  \\
  &+
  \sum_{n=0}^{N}
  \norm{\partial_t^{n+1}u}_{L^2(0,T;H^{M-1})}^{2}.
\end{aligned}
\end{equation}
When \(M=0\), interpret the lower row in the inhomogeneous Bessel-potential
space \(H^{-1}\).  The record \(\cR_{N,M}[u_0;T]\) retains these velocity norms,
their normalized pressure rows, and the equations that identify all entries as
coordinates of \((u,p)\).  Smoothness pays every finite record once \(T<T_*\)
has been fixed.

\begin{lemma}[Rectangles on strict classical subintervals]\label{lem:StrictSubintervalRectangles}
For each strict cutoff \(T<T_*\) and each finite pair \(N,M\),
\begin{equation}\label{eq:AppendixStrictSubintervalRectangleBound}
  \mathfrak R_{N,M}(u;T)<\infty.
\end{equation}
If \(N'\leq N\), \(M'\leq M\), and \(T'\leq T\), forgetting the excess
coordinates and restricting time produces the smaller record.
\end{lemma}

\begin{proof}
The closed interval \([0,T]\) stays strictly inside the classical lifespan.
Persistence puts every fixed \(\partial_t^nu\) continuously in every fixed
\(H^q\) there, and compactness bounds that coordinate.  Hence
\[
  \norm{\partial_t^n u}_{L^2(0,T;H^q)}^2
  \leq
  T\sup_{0\leq t\leq T}
  \norm{\partial_t^n u(t)}_{H^q}^{2}
  <\infty.
\]
Apply this inequality to the finitely many orders in
\eqref{eq:SectionThreeRectangleNorm}.  Restriction changes neither the solution
nor its derivatives, which also proves compatibility.
\end{proof}

Strict-cutoff finiteness is not terminal control.  The scalar function
\(f(t)=(T_*-t)^{-1/2}\) already separates the two quantifiers:
\[
  f\in L^2(0,T)\quad\text{for every }T<T_*,
  \qquad
  f\notin L^2(I_*).
\]
This is not a fluid model; it only shows that every strict interval may be
finite while the terminal integral diverges.

The preceding lemma permits exactly that growth.  The producer theorem
\ref{thm:BodyWholeTerminalRectangleTheorem} supplies the missing fact: after a
finite \((N,M)\) is chosen, its bound is independent of the cutoff.  For
\(N'\leq N\) and \(M'\leq M\), define the terminal bonding map
\[
  \rho_{N',M';I_*}^{N,M;I_*}:
  \mathscr X_{N,M}(I_*)
  \longrightarrow
  \mathscr X_{N',M'}(I_*)
\]
by dropping the higher rows, embedding the retained velocity and pressure
coordinates into the lower spaces, and recomputing the adjacent-row norm on
\(I_*\).  Direct restriction and restriction in stages agree, while an unchanged
pair gives the identity.

\Needspace{22\baselineskip}
\begin{proposition}[Terminal realization of the rectangle estimate]\label[proposition]{thm:WholeTerminalCompatibleRectangleEstimate}
Assume Theorem~\ref{thm:BodyWholeTerminalRectangleTheorem}.  For every finite
\(N,M\), the datum-generated pressure-complete history satisfies
\begin{equation}\label{eq:WholeTerminalRectangleBound}
\begin{aligned}
  \mathfrak R_{N,M}(u;I_*)
  :={}&
  \sum_{n=0}^{N}
  \norm{\partial_t^n u}_{L^2(I_*;H^{M+1})}^{2}
  \\
  &+
  \sum_{n=0}^{N}
  \norm{\partial_t^{n+1}u}_{L^2(I_*;H^{M-1})}^{2}
  <\infty.
\end{aligned}
\end{equation}
The same coordinates therefore define a record on all of \(I_*\):
\[
  \cR_{N,M}[u_0]
  \in
  \mathscr X_{N,M}(I_*),
\]
For every lower pair \((N',M')\), its terminal restriction satisfies
\begin{equation}\label{eq:WholeTerminalRectangleCompatibility}
  \rho_{N',M';I_*}^{N,M;I_*}
  \bigl(\cR_{N,M}[u_0]\bigr)
  =\cR_{N',M'}[u_0].
\end{equation}
All shared entries are the same distributional derivatives of \(u\), and all
pressure entries retain the normalization
\eqref{eq:SectionThreePressureNormalization}.
\end{proposition}

\begin{proof}
Fix the pair \((N,M)\) and, for \(0\leq n\leq N\), set
\[
  f_n^+(t):=\norm{V_n(t)}_{H^{M+1}}^2,
  \qquad
  f_n^-(t):=\norm{V_{n+1}(t)}_{H^{M-1}}^2.
\]
The functions are nonnegative and \((0,T)\) increases to \(I_*\).  Monotone
convergence gives, for either sign,
\[
  \int_{I_*}f_n^\sigma(t)\,dt
  =
  \lim_{T\uparrow T_*}\int_0^T f_n^\sigma(t)\,dt.
\]
Each strict integral is one nonnegative term in the producer bound.  Its limit
is therefore finite, and
\[
  V_n\in L^2(I_*;H^{M+1}),
  \qquad
  V_{n+1}\in L^2(I_*;H^{M-1}).
\]
Summing the finitely many limits gives
\eqref{eq:WholeTerminalRectangleBound}.  The Riesz formula, Sobolev
multiplication, and Cauchy--Schwarz in time also give
\[
  \norm{Q_n}_{L^1(I_*;H^M)}
  \leq
  C_{M,n}\sum_{a=0}^{n}
  \norm{V_a}_{L^2(I_*;H^{M+1})}
  \norm{V_{n-a}}_{L^2(I_*;H^{M+1})}
  <\infty.
\]
Thus the pressure rows belong to the terminal record.  Dropping higher rows or
lowering the spatial height changes no shared derivative or normalized
pressure, proving \eqref{eq:WholeTerminalRectangleCompatibility}.
\end{proof}

For example, \((N,M)=(2,3)\) retains \(V_0,V_1,V_2\) in
\(L^2(I_*;H^4)\), \(V_1,V_2,V_3\) in \(L^2(I_*;H^2)\), and
\(Q_0,Q_1,Q_2\) in \(L^1(I_*;H^3)\).  Restriction to \((1,2)\) drops
\(V_3,Q_2\) and reads the remaining coordinates in the corresponding lower
spaces.  Nothing is recomputed: each retained distribution is the same one.
Compatibility concerns this identity; terminal boundedness was the separate
quantifier paid by the producer theorem.

Each finite coordinate can now be read through its adjacent pair.
\begin{proposition}[Adjacent-row recovery]\label{prop:AdjacentRowRecovery}
For every finite \(n,m\), the producer estimate gives
\begin{equation}\label{eq:WholeTerminalAdjacentRows}
  V_n\in L^2(I_*;H^{m+1}),
  \qquad
  \partial_tV_n=V_{n+1}\in L^2(I_*;H^{m-1}).
\end{equation}
The bound may depend on the chosen pair.
\end{proposition}

\begin{proof}
Choose \((N,M)=(n,m)\).  The two terms indexed by \(n\) in
\eqref{eq:WholeTerminalRectangleBound} are precisely the displayed adjacent
rows.  They are actual derivatives of \(u\), and restriction preserves that
identity.
\end{proof}

The strict and terminal records have always referred to the same derivatives of
\((u,p)\).  What has changed is the interval on which their norms are finite;
\eqref{eq:WholeTerminalRectangleCompatibility} now makes the terminal family
compatible under every decrease of either cutoff.
